\textbf{Phần 1. Trắc nghiệm lựa chọn}
\setcounterpageref{ex}{0}

\begin{ex} %[5X11V63]
Cho bảng số liệu về số học sinh giỏi của $4$ lớp khối 5 như sau: Lớp 5A có $18$ bạn, lớp 5B có $22$ bạn, lớp 5C có $20$ bạn, lớp 5D có $16$ bạn. Trung bình mỗi lớp có bao nhiêu học sinh giỏi?
\choice
{$18\text{ bạn}$}
{\True $19\text{ bạn}$}
{$20\text{ bạn}$}
{$21\text{ bạn}$}
\loigiai{
Tổng số học sinh giỏi của cả $4$ lớp là:
$$18 + 22 + 20 + 16 = 76\text{ (bạn)}.$$
Trung bình mỗi lớp có số học sinh giỏi là:
$$76 : 4 = 19\text{ (bạn)}.$$
\textbf{Đáp án của câu hỏi là B}
}
\end{ex}

\begin{ex} %[5X11V63]
Khảo sát số giờ sử dụng điện thoại thông minh mỗi ngày của một nhóm học sinh thu được dãy số liệu (đơn vị: giờ): $1{,}5;\; 2;\; 1;\; 2{,}5;\; 2;\; 1{,}5;\; 3;\; 2;\; 0{,}5;\; 2$. Giá trị xuất hiện nhiều lần nhất trong dãy số liệu trên là
\choice
{$1{,}5\text{ giờ}$}
{\True $2\text{ giờ}$}
{$2{,}5\text{ giờ}$}
{$1\text{ giờ}$}
\loigiai{
Đếm số lần xuất hiện của các giá trị:
- Giá trị $0{,}5$: xuất hiện $1$ lần.
- Giá trị $1$: xuất hiện $1$ lần.
- Giá trị $1{,}5$: xuất hiện $2$ lần.
- Giá trị $2$: xuất hiện $4$ lần.
- Giá trị $2{,}5$: xuất hiện $1$ lần.
- Giá trị $3$: xuất hiện $1$ lần.
Vậy giá trị $2\text{ giờ}$ xuất hiện nhiều nhất ($4$ lần).\\
\textbf{Đáp án của câu hỏi là B}
}
\end{ex}

\begin{ex} %[5X11V63]
Điểm kiểm tra môn Toán cuối kì I của tổ 1 gồm $8$ bạn là: $9, 8, 10, 7, 9, 8, 9, 8$. Điểm trung bình môn Toán của tổ 1 là
\choice
{$8{,}25$}
{\True $8{,}5$}
{$8{,}75$}
{$8{,}0$}
\loigiai{
Tổng số điểm của $8$ bạn là:
$$9 + 8 + 10 + 7 + 9 + 8 + 9 + 8 = 68.$$
Điểm trung bình của tổ 1 là:
$$68 : 8 = 8{,}5.$$
\textbf{Đáp án của câu hỏi là B}
}
\end{ex}

\begin{ex} %[5X11V63]
Một nhóm từ thiện quyên góp sách giáo khoa cũ trong $5$ ngày. Số lượng sách quyên góp được lần lượt là: $45, 60, 55, 70, 80$ (quyển). Biết mục tiêu của nhóm là đạt trung bình $65$ quyển mỗi ngày. Hỏi nhóm đó có hoàn thành mục tiêu không và vượt hay thiếu bao nhiêu quyển?
\choice
{Thiếu $15$ quyển}
{\True Thiếu $15$ quyển}
{Vượt $15$ quyển}
{Đạt vừa đúng mục tiêu}
\loigiai{
Tổng số sách quyên góp được trong $5$ ngày là:
$$45 + 60 + 55 + 70 + 80 = 310\text{ (quyển)}.$$
Tổng số sách theo mục tiêu đặt ra là:
$$65 \times 5 = 325\text{ (quyển)}.$$
Vì $310 < 325$ nên nhóm chưa hoàn thành mục tiêu và còn thiếu:
$$325 - 310 = 15\text{ (quyển)}.$$
\textbf{Đáp án của câu hỏi là B}
}
\end{ex}

\begin{ex} %[5X11V63]
Cho dãy số liệu về nhiệt độ đo lúc $12\text{ giờ}$ trưa trong một tuần tại một địa phương (đơn vị: $^{\circ}\text{C}$): $28, 30, 32, 31, 33, 29, 31$. Hiệu giữa nhiệt độ cao nhất và nhiệt độ thấp nhất trong tuần đó là
\choice
{$4^{\circ}\text{C}$}
{\True $5^{\circ}\text{C}$}
{$6^{\circ}\text{C}$}
{$7^{\circ}\text{C}$}
\loigiai{
Nhiệt độ cao nhất trong tuần là $33^{\circ}\text{C}$.\\
Nhiệt độ thấp nhất trong tuần là $28^{\circ}\text{C}$.\\
Hiệu nhiệt độ là:
$$33 - 28 = 5^{\circ}\text{C}.$$
\textbf{Đáp án của câu hỏi là B}
}
\end{ex}

\begin{ex} %[5X11V63]
Khảo sát phương tiện đến trường của $40$ học sinh lớp 5A thu được kết quả: $16$ bạn đi xe đạp, $14$ bạn được bố mẹ chở bằng xe máy, $6$ bạn đi bộ và số bạn còn lại đi xe buýt. Tỉ số phần trăm của số học sinh đi xe buýt so với cả lớp là
\choice
{$8\%$}
{\True $10\%$}
{$12\%$}
{$15\%$}
\loigiai{
Số học sinh đi xe buýt là:
$$40 - (16 + 14 + 6) = 4\text{ (bạn)}.$$
Tỉ số phần trăm học sinh đi xe buýt là:
$$4 : 40 \times 100\% = 10\%.$$
\textbf{Đáp án của câu hỏi là B}
}
\end{ex}

\begin{ex} %[5X11V63]
Số lượng cây hoa trồng được trong vườn trường của ba khối 3, 4, 5 được ghi lại: Khối 3 trồng được $120$ cây, Khối 4 trồng được nhiều hơn Khối 3 là $30$ cây, Khối 5 trồng gấp rưỡi Khối 4. Tổng số cây hoa cả ba khối trồng được là
\choice
{$450\text{ cây}$}
{\True $495\text{ cây}$}
{$520\text{ cây}$}
{$480\text{ cây}$}
\loigiai{
Khối 4 trồng được:
$$120 + 30 = 150\text{ (cây)}.$$
Khối 5 trồng được:
$$150 \times 1{,}5 = 225\text{ (cây)}.$$
Tổng số cây cả ba khối trồng được là:
$$120 + 150 + 225 = 495\text{ (cây)}.$$
\textbf{Đáp án của câu hỏi là B}
}
\end{ex}

\begin{ex} %[5X11V63]
Khảo sát số lượng bài tập hoàn thành trong tuần của $10$ học sinh ta có bảng số liệu:
- $4$ bạn làm được $20$ bài;
- $3$ bạn làm được $25$ bài;
- $2$ bạn làm được $30$ bài;
- $1$ bạn làm được $35$ bài.
Tổng số bài tập $10$ bạn đó đã làm là
\choice
{$240\text{ bài}$}
{\True $250\text{ bài}$}
{$260\text{ bài}$}
{$275\text{ bài}$}
\loigiai{
Tổng số bài tập làm được là:
$$(4 \times 20) + (3 \times 25) + (2 \times 30) + (1 \times 35) = 80 + 75 + 60 + 35 = 250\text{ (bài)}.$$
\textbf{Đáp án của câu hỏi là B}
}
\end{ex}

\begin{ex} %[5X11V63]
Trong một đợt kiểm tra sức khỏe, chiều cao của $5$ bạn trong đội tuyển bóng rổ trường lần lượt là: $152\text{ cm}, 148\text{ cm}, 156\text{ cm}, 154\text{ cm}, 150\text{ cm}$. Nếu bổ sung thêm bạn Minh cao $162\text{ cm}$ vào đội tuyển thì chiều cao trung bình của đội thay đổi thế nào?
\choice
{Tăng thêm $1\text{ cm}$}
{\True Tăng thêm $1{,}67\text{ cm}$}
{Tăng thêm $2\text{ cm}$}
{Không đổi}
\loigiai{
Chiều cao trung bình ban đầu của 5 bạn:
$$(152 + 148 + 156 + 154 + 150) : 5 = 760 : 5 = 152\text{ (cm)}.$$
Khi thêm bạn Minh, tổng chiều cao của 6 bạn là:
$$760 + 162 = 922\text{ (cm)}.$$
Chiều cao trung bình mới là:
$$922 : 6 \approx 153{,}67\text{ (cm)}.$$
Chiều cao trung bình tăng thêm:
$$153{,}67 - 152 = 1{,}67\text{ (cm)}.$$
\textbf{Đáp án của câu hỏi là B}
}
\end{ex}

\begin{ex} %[5X11V63]
Số lượng sữa chua bán ra trong $4$ ngày đầu tuần của một cửa hàng là: Thứ Hai: $45$ hộp; Thứ Ba: $50$ hộp; Thứ Tư: $65$ hộp; Thứ Năm: $60$ hộp. Để trung bình mỗi ngày trong cả tuần ($7$ ngày) bán được $60$ hộp thì trong $3$ ngày cuối tuần cửa hàng phải bán được tất cả bao nhiêu hộp sữa chua?
\choice
{$180\text{ hộp}$}
{\True $200\text{ hộp}$}
{$210\text{ hộp}$}
{$220\text{ hộp}$}
\loigiai{
Tổng số hộp sữa chua cần bán trong cả $7$ ngày là:
$$60 \times 7 = 420\text{ (hộp)}.$$
Số hộp sữa chua đã bán trong $4$ ngày đầu tuần là:
$$45 + 50 + 65 + 60 = 220\text{ (hộp)}.$$
Số hộp sữa chua cần bán trong $3$ ngày cuối tuần là:
$$420 - 220 = 200\text{ (hộp)}.$$
\textbf{Đáp án của câu hỏi là B}
}
\end{ex}

\begin{ex} %[5X11V63]
Dãy số liệu nào dưới đây đã được sắp xếp theo thứ tự giảm dần?
\choice
{$12{,}5;\; 12{,}05;\; 12{,}52;\; 11{,}9$}
{\True $15{,}4;\; 15{,}35;\; 15{,}08;\; 14{,}9$}
{$8{,}7;\; 8{,}65;\; 8{,}68;\; 8{,}5$}
{$20{,}1;\; 20{,}01;\; 20{,}2;\; 19{,}9$}
\loigiai{
Xét dãy B: $15{,}4 > 15{,}35 > 15{,}08 > 14{,}9$ là dãy số giảm dần chính xác.\\
\textbf{Đáp án của câu hỏi là B}
}
\end{ex}

\begin{ex} %[5X11V63]
Một nhóm khảo sát hỏi ý kiến $50$ người về việc thích đọc sách, kết quả thu được: $28$ người rất thích, $16$ người thích vừa phải và số người còn lại không thích. Tỉ số phần trăm của số người không thích đọc sách là
\choice
{$10\%$}
{\True $12\%$}
{$14\%$}
{$8\%$}
\loigiai{
Số người không thích đọc sách là:
$$50 - (28 + 16) = 6\text{ (người)}.$$
Tỉ số phần trăm là:
$$6 : 50 \times 100\% = 12\%.$$
\textbf{Đáp án của câu hỏi là B}
}
\end{ex}

\textbf{Phần 2. Trắc nghiệm đúng sai}
\setcounterpageref{ex}{0}

\begin{ex} %[5X11V63]
Bảng thống kê số lượng điểm $9$ và điểm $10$ của $4$ bạn trong tháng như sau: An ($8$ điểm), Bình ($10$ điểm), Cúc ($6$ điểm), Dũng ($12$ điểm). Xét tính đúng/sai của các nhận định sau:
\choiceTF
{\True Tổng số điểm $9$ và điểm $10$ của cả $4$ bạn là $36$ điểm}
{\True Trung bình mỗi bạn đạt được $9$ điểm tốt ($9$ hoặc $10$)}
{Bạn An có số điểm tốt bằng một nửa bạn Dũng}
{\True Số điểm tốt của bạn Dũng gấp đôi số điểm tốt của bạn Cúc}
\loigiai{
\begin{itemchoice}
\item Đúng. $8 + 10 + 6 + 12 = 36$ điểm.
\item Đúng. $36 : 4 = 9$ điểm.
\item Sai. An có 8 điểm, Dũng có 12 điểm, tỉ số là $8 : 12 = \dfrac{2}{3}$ chứ không phải một nửa.
\item Đúng. Dũng có 12 điểm, Cúc có 6 điểm $\Rightarrow 12 : 6 = 2$ (gấp đôi).
\end{itemchoice}
}
\end{ex}

\begin{ex} %[5X11V63]
Khảo sát số ki-lô-gam giấy vụn thu gom của $3$ lớp: Lớp 5A thu được $120\text{ kg}$, lớp 5B thu được bằng $\dfrac{5}{6}$ lớp 5A, lớp 5C thu được nhiều hơn lớp 5B là $15\text{ kg}$. Xét tính đúng/sai của các nhận định sau:
\choiceTF
{\True Lớp 5B thu gom được $100\text{ kg}$ giấy vụn}
{\True Lớp 5C thu gom được $115\text{ kg}$ giấy vụn}
{Lớp 5A thu gom được ít giấy vụn nhất trong ba lớp}
{\True Cả ba lớp thu gom được tất cả $335\text{ kg}$ giấy vụn}
\loigiai{
\begin{itemchoice}
\item Đúng. $120 \times \dfrac{5}{6} = 100\text{ kg}$.
\item Đúng. $100 + 15 = 115\text{ kg}$.
\item Sai. Lớp 5A thu được $120\text{ kg}$ là nhiều nhất chứ không phải ít nhất.
\item Đúng. $120 + 100 + 115 = 335\text{ kg}$.
\end{itemchoice}
}
\end{ex}

\textbf{Phần 3. Câu trả lời ngắn}
\setcounterpageref{ex}{0}

\begin{ex} %[5X11V63]
Cho dãy số liệu về số học sinh của $5$ lớp khối 5: $32, 34, 30, 36, 38$. Hỏi trung bình mỗi lớp có bao nhiêu học sinh?
\shortans{34}
\loigiai{
Tổng số học sinh của $5$ lớp là:
$$32 + 34 + 30 + 36 + 38 = 170\text{ (học sinh)}.$$
Trung bình mỗi lớp có số học sinh là:
$$170 : 5 = 34\text{ (học sinh)}.$$
}
\end{ex}

\begin{ex} %[5X11V63]
Số lượng sách đọc được của một bạn trong $4$ tháng đầu năm là: Tháng 1: $6$ cuốn; Tháng 2: $4$ cuốn; Tháng 3: $8$ cuốn; Tháng 4: $10$ cuốn. Để trung bình mỗi tháng trong $5$ tháng bạn đọc được $8$ cuốn thì Tháng 5 bạn phải đọc bao nhiêu cuốn sách?
\shortans{12}
\loigiai{
Tổng số sách cần đọc trong $5$ tháng là:
$$8 \times 5 = 40\text{ (cuốn)}.$$
Số sách đã đọc trong $4$ tháng đầu là:
$$6 + 4 + 8 + 10 = 28\text{ (cuốn)}.$$
Số sách cần đọc trong Tháng 5 là:
$$40 - 28 = 12\text{ (cuốn)}.$$
}
\end{ex}

\begin{ex} %[5X11V63]
Một cửa hàng thống kê số bao gạo bán được trong $3$ ngày: Ngày thứ nhất bán $45$ bao, ngày thứ hai bán $55$ bao, ngày thứ ba bán được số bao gạo bằng trung bình cộng của cả $3$ ngày. Hỏi ngày thứ ba cửa hàng bán được bao nhiêu bao gạo?
\shortans{50}
\loigiai{
Vì số bao gạo ngày thứ ba bằng trung bình cộng của cả $3$ ngày nên ngày thứ ba bán đúng bằng trung bình cộng của hai ngày đầu:\\
$$(45 + 55) : 2 = 50\text{ (bao)}.$$
}
\end{ex}

\begin{ex} %[5X11V63]
Trong đợt thi đua chào mừng ngày Nhà giáo Việt Nam, tổ Một đạt được $120$ hoa điểm tốt, tổ Hai đạt được $135$ hoa điểm tốt, tổ Ba đạt được nhiều hơn mức trung bình của cả ba tổ là $15$ hoa điểm tốt. Hỏi tổ Ba đạt được bao nhiêu hoa điểm tốt?
\shortans{150}
\loigiai{
Hai lần mức trung bình cộng của cả $3$ tổ là:
$$120 + 135 + 15 = 270\text{ (hoa)}.$$
Mức trung bình cộng của cả $3$ tổ là:
$$270 : 2 = 135\text{ (hoa)}.$$
Số hoa điểm tốt của tổ Ba là:
$$135 + 15 = 150\text{ (hoa)}.$$
}
\end{ex}

\begin{ex} %[5X11V63]
Một gia đình ghi chép số số điện tiêu thụ (theo $\text{kWh}$) trong $6$ tháng đầu năm: $140, 160, 180, 220, 260, 240$. Hỏi số điện tiêu thụ nhiều nhất hơn số điện tiêu thụ ít nhất trong $6$ tháng đó là bao nhiêu $\text{kWh}$?
\shortans{120}
\loigiai{
Số điện tiêu thụ nhiều nhất là $260\text{ kWh}$.\\
Số điện tiêu thụ ít nhất là $140\text{ kWh}$.\\
Hiệu số điện là:
$$260 - 140 = 120\text{ (kWh)}.$$
}
\end{ex}

\begin{ex} %[5X11V63]
Một lớp học có $35$ học sinh. Sau khi khảo sát môn thể thao yêu thích nhất thì có $40\%$ số bạn thích bóng đá, $20\%$ số bạn thích cầu lông, số còn lại thích bơi lội. Hỏi có bao nhiêu bạn thích bơi lội?
\shortans{14}
\loigiai{
Tỉ số phần trăm các bạn thích bơi lội là:
$$100\% - (40\% + 20\%) = 40\%.$$
Số bạn thích bơi lội là:
$$35 \times 40\% = 14\text{ (bạn)}.$$
}
\end{ex}
